%Introduction---------------

\setlength{\parindent}{2em} {在现代交通运输体系快速发展的背景下，高速、安全、绿色、低噪声已成为轨道交通技术升级的重要方向。传统轮轨列车受轮轨黏着、机械磨耗和运行噪声等因素限制，在进一步提升速度和降低维护成本方面面临一定瓶颈。磁浮列车通过电磁力实现车体与轨道之间的非接触支承，能够有效减少机械摩擦和部件磨损，具有运行速度高、乘坐舒适性好、维护成本低等优势，是未来高速地面交通的重要发展方向之一。

抱轨式磁浮列车作为高速磁浮交通系统的重要形式，其稳定运行依赖于悬浮电磁铁、悬浮架、空气弹簧以及车体之间的协调作用。列车运行过程中，电磁铁产生的吸引力需要实时平衡车体重力及外界扰动，使悬浮间隙保持在合理范围内。然而，电磁力与电流、悬浮间隙之间存在明显的非线性关系；同时，车辆在高速运行时还会受到轨道不平顺、气流扰动、电磁干扰和温度变化等因素影响，这些因素都可能引起悬浮间隙波动，甚至诱发悬浮系统故障。

在磁浮列车悬浮控制系统中，电磁铁功率放大器直接影响电磁力的实际输出。当功率放大器出现异常时，电磁铁产生的实际电磁力可能相较理想值发生突增或衰减，进而破坏车体、悬浮架与轨道之间的受力平衡。若不能及时识别此类故障，不仅会降低列车运行平稳性，还可能造成悬浮间隙异常变化，影响系统安全。因此，如何利用电流、悬浮间隙、车体加速度和实际电磁力等监测数据，建立可靠的电磁力模型和垂向动力学模型，并进一步设计功率放大器故障检测方法，具有重要的工程价值和实际意义。}
